\begin{titlepage}
\centering
\setlength{\parindent}{0pt}

\hbox{}
\vspace{2cm}

{\fontsize{36}{36}\selectfont\textbf{Զ\;Ր\;ՈՒ\;Յ\;Ց\;Ն\;Ե\;Ր}}
\null\vskip12pt\null
{\fontsize{64}{64}\selectfont\textbf{Ս\;Ի}}
\null\vskip12pt\null
{\fontsize{28}{28}\selectfont\textbf{Լ\,Ե\,Զ\,Վ\,Ի\; Մ\,Ա\,Ս\,Ի\,Ն}}

\vspace{0.6cm}
\rule{0.3\textwidth}{0.4pt}\par
\vspace{0.6cm}
{\Large \textit{Արմեն~~Բադալյան}}

\vfill

ԱՐՄԵՆԻԱ • ԵՐԵՒԱՆ

\MakeUppercase{\romannumeral\year{}}

\end{titlepage}


\begin{titlepage}

\begingroup
\small\sffamily\linespread{1}
\begin{tabular}{rl}
ԳՄԴ & xxx.xx:xxx.xx \\
ՀՏԴ & xx.xx \\
Ա & xx \\
\end{tabular}
\endgroup

\vskip2cm

\begin{tabular}{rl}
     & Բադալյան Արմեն \\
Բ xx & \textsl{Զրույցներ Սի լեզվի մասին} / Արմեն Բադալյան \\
     & Երևան, \the\year, xxx էջ
\end{tabular}

\begin{small}
Սա գիրք է Սի ծրագրավորման լեզվի հիմունքների մասին։ Հեղինակը պարզ լեզվով
և պարզ օրինակներով պատմում է, թե ինչպես կարելի է սկսել ծրագրավորել Սի
լեզվով։ Բացի լեզվական հիմնական կառուցվածքները, դիտարկվում են նաև այնպիսի
թեմաներ, ինչպիսիք են ստատիկ ու դինամիկ գրադարանների կառուցումը, ֆայլերի
ու տեքստերի հետ աշխատանքը, հոսքերը։ Գիրքը նախատեսված է ավագ դասարանների
դպրոցականների, ուսանողների և բոլոր հետաքրքրվողների համար։
\end{small}


\vfill
\textsf{\copyright\ Արմեն Բադալյան, \the\year։}

\end{titlepage}


\begin{titlepage}
\hbox{}
\vskip5cm
\begin{center}
\textit{....ին։}
\end{center}
\vfill
\end{titlepage}
